% TOP_PROBLEM: {"problemId": "problem.tates-rational-leading-term-stark-conjecture", "canonicalTitle": "Tate's rational leading-term Stark conjecture", "releaseRank": 166, "primaryDomain": "number_theory_arithmetic_geometry", "primaryDomainLabel": "Number theory & arithmetic geometry", "displayStatus": "open_with_solved_subcases", "releaseStatus": "open_with_solved_subcases"}
% !TeX program = lualatex
% Definition and sources only; this file is not an LLM solution attempt.
\documentclass[11pt]{article}
\usepackage[margin=25mm]{geometry}
\usepackage{fontspec}
\setmainfont{Latin Modern Roman}
\usepackage{amsmath,amssymb,mathtools}
\usepackage{unicode-math}
\setmathfont{Latin Modern Math}
\usepackage{xurl,hyperref}
\hypersetup{colorlinks=true,urlcolor=blue}
\setcounter{secnumdepth}{0}
\setlength{\parindent}{0pt}
\setlength{\parskip}{5pt}
\setlength{\emergencystretch}{3em}
\begin{document}
\section*{\texorpdfstring{Tate's rational leading-term Stark conjecture}{Tate's rational leading-term Stark conjecture}}\label{166}
\subsection{Definitions and mathematical statement}
For a Galois extension $L/K$ with group $G$, let $S_L$ be the places above $S$ and set
$X=\ker({\mathbb{Q}}[S_L]\xrightarrow{\sum}{\mathbb{Q}})$ and
$U={\mathbb{Q}}\otimes_{{\mathbb{Z}}}\mathcal O_{L,S}^{\times}$. For a ${\mathbb{Q}}[G]$-isomorphism $f:X\to U$, the logarithmic map
$\lambda(u)=-\sum_{w\in S_L}\log|u|_w\,w$
uses product-formula normalized absolute values. For an irreducible character $\chi$, take the determinant of $\lambda f$ on
$\operatorname{Hom}_{{\mathbb{C}}[G]}(V_{\chi^\vee},{\mathbb{C}}\otimes X)$ to obtain $R_S(\chi,f)$. Let $L_S^*(0,\chi)$ be the first nonzero Laurent coefficient at zero of the Artin $L$-function with the $S$ factors omitted. The target is
\[
 A_S(\chi,f):=\frac{R_S(\chi,f)}{L_S^*(0,\chi)}\in{\mathbb{Q}}(\chi),
 \qquad\sigma A_S(\chi,f)=A_S(\chi^\sigma,f).
\]
The empty determinant is one; the character dual and the sign in $\lambda$ are part of the chosen normalization.
\par\smallskip\textit{Formulation and concept references:} \href{https://arxiv.org/pdf/2106.05619}{[S1]}.

\subsection{Short English statement}
After dividing the appropriate logarithmic regulator by an Artin L-function's leading term, is the result algebraic in the character field and compatible with every field automorphism?
\subsection{Sources}
\begin{itemize}
\item[S1]Andreas Nickel, The strong Stark conjecture for totally odd characters; arXiv2106.05619v1 June10,2021; live metadata lists ProcAMS152(2024)147–162.
\url{https://arxiv.org/pdf/2106.05619}.
\textit{Formulation source.}
\end{itemize}
\end{document}
